% Options for packages loaded elsewhere
\PassOptionsToPackage{unicode}{hyperref}
\PassOptionsToPackage{hyphens}{url}
%
\documentclass[
  11pt,
]{article}
\usepackage{amsmath,amssymb}
\usepackage{iftex}
\ifPDFTeX
  \usepackage[T1]{fontenc}
  \usepackage[utf8]{inputenc}
  \usepackage{textcomp} % provide euro and other symbols
\else % if luatex or xetex
  \usepackage{unicode-math} % this also loads fontspec
  \defaultfontfeatures{Scale=MatchLowercase}
  \defaultfontfeatures[\rmfamily]{Ligatures=TeX,Scale=1}
\fi
\usepackage{lmodern}
\ifPDFTeX\else
  % xetex/luatex font selection
\fi
% Use upquote if available, for straight quotes in verbatim environments
\IfFileExists{upquote.sty}{\usepackage{upquote}}{}
\IfFileExists{microtype.sty}{% use microtype if available
  \usepackage[]{microtype}
  \UseMicrotypeSet[protrusion]{basicmath} % disable protrusion for tt fonts
}{}
\makeatletter
\@ifundefined{KOMAClassName}{% if non-KOMA class
  \IfFileExists{parskip.sty}{%
    \usepackage{parskip}
  }{% else
    \setlength{\parindent}{0pt}
    \setlength{\parskip}{6pt plus 2pt minus 1pt}}
}{% if KOMA class
  \KOMAoptions{parskip=half}}
\makeatother
\usepackage{xcolor}
\usepackage[margin=1in]{geometry}
\usepackage{longtable,booktabs,array}
\usepackage{calc} % for calculating minipage widths
% Correct order of tables after \paragraph or \subparagraph
\usepackage{etoolbox}
\makeatletter
\patchcmd\longtable{\par}{\if@noskipsec\mbox{}\fi\par}{}{}
\makeatother
% Allow footnotes in longtable head/foot
\IfFileExists{footnotehyper.sty}{\usepackage{footnotehyper}}{\usepackage{footnote}}
\makesavenoteenv{longtable}
\usepackage{graphicx}
\makeatletter
\def\maxwidth{\ifdim\Gin@nat@width>\linewidth\linewidth\else\Gin@nat@width\fi}
\def\maxheight{\ifdim\Gin@nat@height>\textheight\textheight\else\Gin@nat@height\fi}
\makeatother
% Scale images if necessary, so that they will not overflow the page
% margins by default, and it is still possible to overwrite the defaults
% using explicit options in \includegraphics[width, height, ...]{}
\setkeys{Gin}{width=\maxwidth,height=\maxheight,keepaspectratio}
% Set default figure placement to htbp
\makeatletter
\def\fps@figure{htbp}
\makeatother
\setlength{\emergencystretch}{3em} % prevent overfull lines
\providecommand{\tightlist}{%
  \setlength{\itemsep}{0pt}\setlength{\parskip}{0pt}}
\setcounter{secnumdepth}{-\maxdimen} % remove section numbering
\ifLuaTeX
\usepackage[bidi=basic]{babel}
\else
\usepackage[bidi=default]{babel}
\fi
\babelprovide[main,import]{spanish}
% get rid of language-specific shorthands (see #6817):
\let\LanguageShortHands\languageshorthands
\def\languageshorthands#1{}
\ifLuaTeX
  \usepackage{selnolig}  % disable illegal ligatures
\fi
\IfFileExists{bookmark.sty}{\usepackage{bookmark}}{\usepackage{hyperref}}
\IfFileExists{xurl.sty}{\usepackage{xurl}}{} % add URL line breaks if available
\urlstyle{same}
\hypersetup{
  pdftitle={Estadística para la Toma de Decisiones},
  pdfauthor={Alexander Buriticá Casanova},
  pdflang={es},
  hidelinks,
  pdfcreator={LaTeX via pandoc}}

\title{Estadística para la Toma de Decisiones}
\usepackage{etoolbox}
\makeatletter
\providecommand{\subtitle}[1]{% add subtitle to \maketitle
  \apptocmd{\@title}{\par {\large #1 \par}}{}{}
}
\makeatother
\subtitle{Programa del curso -- Pregrado}
\author{Alexander Buriticá Casanova}
\date{Periodo académico 2026-1}

\begin{document}
\maketitle

\hypertarget{informaciuxf3n-general-del-curso}{%
\subsection{Información general del
curso}\label{informaciuxf3n-general-del-curso}}

\textbf{Curso:} BA2 -- Estadística para la Toma de Decisiones\\
\textbf{Curso:} ECO-06326

\textbf{Programa:} Business Analytics\\
\textbf{Periodo académico:} 2026-1\\
\textbf{Créditos:} 3

\textbf{Profesor:} Alexander Buriticá Casanova, PhD\\
\textbf{Monitora:} María Claudia

\textbf{Día y horario:} Lunes, 14:00-16:59\\
\textbf{Modalidad:} Presencial\\
\textbf{Lugar:} Salón 303, Edificio C

\hypertarget{descripciuxf3n-del-curso}{%
\section{Descripción del curso}\label{descripciuxf3n-del-curso}}

Este curso introduce herramientas fundamentales de estadística
inferencial aplicadas a la toma de decisiones empresariales. El énfasis
está en el uso de datos reales para formular preguntas relevantes de
negocio, analizar información, interpretar resultados y sustentar
decisiones bajo incertidumbre.

El curso combina fundamentos conceptuales con un componente práctico en
\textbf{R}, permitiendo a los estudiantes desarrollar habilidades
analíticas transferibles a contextos profesionales en marketing,
finanzas, economía aplicada y analítica de negocios.

\hypertarget{resultados-de-aprendizaje}{%
\section{Resultados de aprendizaje}\label{resultados-de-aprendizaje}}

Al finalizar el curso, el estudiante será capaz de:

\begin{itemize}
\tightlist
\item
  Formular preguntas de negocio y traducirlas en preguntas estadísticas.
\item
  Describir y explorar datos para identificar patrones relevantes.
\item
  Estimar parámetros e interpretar intervalos de confianza.
\item
  Diseñar y ejecutar pruebas de hipótesis para apoyar decisiones
  empresariales.
\item
  Aplicar ANOVA y regresión lineal simple para responder preguntas
  aplicadas.
\item
  Utilizar \textbf{R} para análisis de datos, visualización e inferencia
  estadística.
\item
  Comunicar hallazgos de manera clara, indicando supuestos, limitaciones
  e implicaciones para la toma de decisiones.
\end{itemize}

\hypertarget{metodologuxeda}{%
\section{Metodología}\label{metodologuxeda}}

El curso combina clases magistrales breves orientadas a la comprensión
conceptual con espacios prácticos de aplicación. Cada sesión se
estructura, en general, en tres momentos:

\begin{enumerate}
\def\labelenumi{\arabic{enumi}.}
\tightlist
\item
  \textbf{Exposición conceptual}: presentación de los fundamentos
  estadísticos, su intuición, supuestos y correcta interpretación.
\item
  \textbf{Taller práctico en clase}: análisis de datos reales utilizando
  \textbf{R}, con énfasis en limpieza de datos, construcción de
  variables e interpretación de resultados.
\item
  \textbf{Actividad corta de verificación}: ejercicio aplicado o quiz
  breve para consolidar los conceptos trabajados.
\end{enumerate}

Adicionalmente, el curso incluye quices periódicos y un proyecto
aplicado progresivo que simula un análisis profesional orientado a la
toma de decisiones empresariales.

\hypertarget{evaluaciuxf3n}{%
\section{Evaluación}\label{evaluaciuxf3n}}

\begin{longtable}[]{@{}lc@{}}
\toprule\noalign{}
Componente & Porcentaje \\
\midrule\noalign{}
\endhead
\bottomrule\noalign{}
\endlastfoot
Talleres prácticos en clase & 20\% \\
Proyecto final (3 entregas) & 25\% \\
Quices (mínimo 4) & 20\% \\
Examen parcial (Parcial 1) & 15\% \\
Examen final acumulativo & 20\% \\
\textbf{Total} & \textbf{100\%} \\
\end{longtable}

\hypertarget{descripciuxf3n-de-los-componentes-de-evaluaciuxf3n}{%
\section{Descripción de los componentes de
evaluación}\label{descripciuxf3n-de-los-componentes-de-evaluaciuxf3n}}

\hypertarget{talleres-pruxe1cticos-en-clase-20}{%
\subsection{Talleres prácticos en clase
(20\%)}\label{talleres-pruxe1cticos-en-clase-20}}

Los talleres constituyen el núcleo aplicado del curso. En ellos, los
estudiantes trabajan con bases de datos reales de marketing, finanzas y
economía aplicada, aplicando las herramientas estadísticas vistas en
clase mediante el uso de \textbf{R}.

Bajo la guía del profesor, los estudiantes limpian datos, construyen
variables, ejecutan análisis inferenciales e interpretan resultados para
respaldar decisiones empresariales. Cada taller incluye una actividad
corta de evaluación asociada al contenido trabajado en la sesión.

\hypertarget{proyecto-final-25}{%
\subsection{Proyecto final (25\%)}\label{proyecto-final-25}}

El proyecto final es el eje integrador del curso y simula el desarrollo
de un análisis profesional orientado a la toma de decisiones
empresariales. Cada grupo debe formular una pregunta de negocio concreta
y responderla utilizando datos reales y las herramientas estadísticas
del curso.

\textbf{Ejemplos de preguntas de proyecto:}

\begin{itemize}
\tightlist
\item
  ¿Una nueva campaña aumenta significativamente la tasa de conversión?
\item
  ¿Reducir el precio incrementa las ventas netas de manera
  estadísticamente significativa?
\item
  ¿Existen diferencias significativas en el comportamiento entre
  segmentos de clientes?
\item
  ¿Qué factores explican el nivel de ventas o el riesgo de
  incumplimiento?
\end{itemize}

\hypertarget{entregas-del-proyecto}{%
\subsubsection{Entregas del proyecto}\label{entregas-del-proyecto}}

\begin{itemize}
\tightlist
\item
  \textbf{Entrega 1 (Semana 5 -- 5\%)}: Contexto del problema,
  formulación de la pregunta de negocio, preguntas estadísticas,
  variables relevantes e hipótesis preliminares.
\item
  \textbf{Entrega 2 (Semana 10 -- 5\%)}: Base de datos depurada y
  análisis descriptivo (tablas y gráficos).
\item
  \textbf{Entrega 3 (01 de junio de 2026 -- 15\%)}: Análisis final,
  interpretación de resultados, recomendaciones y sustentación oral.
\end{itemize}

La sustentación oral es obligatoria. Por limitaciones de tiempo, el
grupo que presenta y la persona expositora podrán definirse de manera
aleatoria.

\hypertarget{quices-20}{%
\subsection{Quices (20\%)}\label{quices-20}}

Se realizarán al menos cuatro quices cortos a lo largo del semestre,
orientados a evaluar la comprensión conceptual y fomentar el estudio
continuo. Los quices podrán ser anunciados o aplicados de manera
sorpresa, y se realizarán en modalidad presencial o virtual según la
dinámica del curso.

\hypertarget{exuxe1menes-35}{%
\subsection{Exámenes (35\%)}\label{exuxe1menes-35}}

\begin{itemize}
\tightlist
\item
  \textbf{Parcial 1 (15\%)}: Evaluación integradora de los contenidos de
  fundamentos y estimación.
\item
  \textbf{Examen final acumulativo (20\%)}: Evaluación de la totalidad
  de los contenidos del curso, con énfasis en interpretación, análisis y
  aplicación.
\end{itemize}

\hypertarget{cronograma-del-curso}{%
\section{Cronograma del curso}\label{cronograma-del-curso}}

Las sesiones del curso se dictan únicamente los \textbf{lunes}. El
cronograma se ajusta a las fechas académicas del periodo 2026-1 y
excluye dos lunes festivos: \textbf{23 de marzo} y \textbf{18 de mayo}.

\begin{longtable}[]{@{}
  >{\raggedright\arraybackslash}p{(\columnwidth - 4\tabcolsep) * \real{0.0702}}
  >{\raggedright\arraybackslash}p{(\columnwidth - 4\tabcolsep) * \real{0.1316}}
  >{\raggedright\arraybackslash}p{(\columnwidth - 4\tabcolsep) * \real{0.7982}}@{}}
\toprule\noalign{}
\begin{minipage}[b]{\linewidth}\raggedright
Sesión
\end{minipage} & \begin{minipage}[b]{\linewidth}\raggedright
Fecha (lunes)
\end{minipage} & \begin{minipage}[b]{\linewidth}\raggedright
Tema
\end{minipage} \\
\midrule\noalign{}
\endhead
\bottomrule\noalign{}
\endlastfoot
1 & 02-feb-2026 & Introducción al curso. Estadística descriptiva. Repaso
básico de R para análisis de datos \\
2 & 09-feb-2026 & Muestreo probabilístico y no probabilístico. Sesgos y
validez \\
3 & 16-feb-2026 & Distribuciones muestrales. Variabilidad y error
estándar \\
4 & 23-feb-2026 & Teorema del Límite Central. Simulaciones y práctica
aplicada en R \\
5 & 02-mar-2026 & Estimación puntual. Incertidumbre y comunicación de
resultados \\
6 & 09-mar-2026 & Intervalos de confianza para la media \\
7 & 16-mar-2026 & Intervalos de confianza para proporciones. Repaso
Bloques I--II \\
--- & 23-mar-2026 & \textbf{Festivo (no hay clase)} \\
8 & 30-mar-2026 & \textbf{Parcial 1} \\
9 & 06-abr-2026 & Pruebas de hipótesis. p-value. Errores tipo I y II \\
10 & 13-abr-2026 & Pruebas de hipótesis para una media y una
proporción \\
11 & 20-abr-2026 & Comparación de dos medias y análisis de varianza
básica \\
12 & 27-abr-2026 & \textbf{Parcial 2} \\
13 & 04-may-2026 & Análisis de varianza (ANOVA). Supuestos e
interpretación \\
14 & 11-may-2026 & Regresión lineal simple. Interpretación aplicada \\
--- & 18-may-2026 & \textbf{Festivo (no hay clase)} \\
15 & 25-may-2026 & Regresión lineal: aplicación en R y taller
integrador \\
16 & 01-jun-2026 & Integración final. \textbf{Entrega 3 y sustentación
del proyecto} \\
\end{longtable}

\hypertarget{reglas-buxe1sicas-del-curso}{%
\section{Reglas básicas del curso}\label{reglas-buxe1sicas-del-curso}}

\begin{itemize}
\tightlist
\item
  Se espera asistencia y participación activa en los talleres prácticos.
\item
  El trabajo en \textbf{R} hace parte central del proceso de aprendizaje
  y será evaluado.
\item
  Las entregas deben presentarse en las fechas establecidas.
\item
  Las políticas de retraso y las rúbricas de evaluación se comunicarán
  oportunamente.
\item
  La honestidad académica aplica a todas las actividades del curso.
\end{itemize}

\end{document}
